\begin{thebibliography}{99}
\bibitem{fefferman}
C. L. Fefferman. \emph{Existence and Smoothness of the Navier--Stokes Equation}. Official problem description, six-page version with errata, Clay Mathematics Institute. User-supplied source; public counterpart accessed 10 September 2026. \url{https://www.claymath.org/wp-content/uploads/2022/06/navierstokes.pdf}.

\bibitem{claystatus}
Clay Mathematics Institute. \emph{Navier--Stokes Equation}. Millennium Prize problem status page, consulted 10 September 2026. \url{https://www.claymath.org/millennium/navier-stokes-equation/}.

\bibitem{openains}
OpenAI. \emph{Finite Time Blowup for Navier--Stokes}. Research manuscript released September 2026, 166 pp.; Theorem 1.1 and introductory exposition consulted. Reported result, not independently verified in the present work. \url{https://cdn.openai.com/pdf/32d9f210-8b73-45e0-91bc-82a30aef8a9a/navier-stokes.pdf}.

\bibitem{openaiannounce}
OpenAI. \emph{On the Navier--Stokes Millennium Prize Problem}. Research announcement, 8 September 2026. \url{https://openai.com/index/navier-stokes-solution/}.

\bibitem{shannon}
C. E. Shannon. A mathematical theory of communication. \emph{Bell System Technical Journal}, 27:379--423 and 623--656, 1948. Author's paper, corrected reprint. \url{https://people.math.harvard.edu/~ctm/home/text/others/shannon/entropy/entropy.pdf}.

\bibitem{nist}
National Institute of Standards and Technology. \emph{Air gap}. Computer Security Resource Center glossary. Consulted 10 September 2026. \url{https://csrc.nist.gov/glossary/term/air_gap}.

\bibitem{noaa}
NOAA Space Weather Prediction Center. \emph{HF Radio Communications}. Official technical explainer, consulted 10 September 2026. \url{https://www.spaceweather.gov/impacts/hf-radio-communications}.

\bibitem{tao}
T. Tao. Finite time blowup for an averaged three-dimensional Navier--Stokes equation. \emph{Journal of the American Mathematical Society}, 29:601--674, 2016. \href{https://arxiv.org/abs/1402.0290}{arXiv:1402.0290}.

\bibitem{cardona3}
R. Cardona, E. Miranda, D. Peralta-Salas, and F. Presas. Constructing Turing complete Euler flows in dimension 3. \emph{Proceedings of the National Academy of Sciences}, 118(19):e2026818118, 2021. \href{https://doi.org/10.1073/pnas.2026818118}{doi:10.1073/pnas.2026818118}. \href{https://arxiv.org/abs/2012.12828}{arXiv:2012.12828}.

\bibitem{cardonaeuclid}
R. Cardona, E. Miranda, and D. Peralta-Salas. Computability and Beltrami fields in Euclidean space. \emph{Journal de Math\'ematiques Pures et Appliqu\'ees}, 169:50--81, 2023. \href{https://doi.org/10.1016/j.matpur.2022.11.007}{doi:10.1016/j.matpur.2022.11.007}. \href{https://arxiv.org/abs/2111.03559}{arXiv:2111.03559}.

\bibitem{dambre}
J. Dambre, D. Verstraeten, B. Schrauwen, and S. Massar. Information processing capacity of dynamical systems. \emph{Scientific Reports}, 2:514, 2012. \href{https://doi.org/10.1038/srep00514}{doi:10.1038/srep00514}.

\bibitem{wright}
L. G. Wright, T. Onodera, M. M. Stein, T. Wang, D. T. Schachter, Z. Hu, and P. L. McMahon. Deep physical neural networks trained with backpropagation. \emph{Nature}, 601:549--555, 2022. \href{https://doi.org/10.1038/s41586-021-04223-6}{doi:10.1038/s41586-021-04223-6}. Preprint under an earlier title: \href{https://arxiv.org/abs/2104.13386}{arXiv:2104.13386}.

\bibitem{sindy}
S. L. Brunton, J. L. Proctor, and J. N. Kutz. Discovering governing equations from data by sparse identification of nonlinear dynamical systems. \emph{Proceedings of the National Academy of Sciences}, 113:3932--3937, 2016. \href{https://doi.org/10.1073/pnas.1517384113}{doi:10.1073/pnas.1517384113}; \href{https://arxiv.org/abs/1509.03580}{arXiv:1509.03580}.

\bibitem{schneider}
F. B. Schneider. Implementing fault-tolerant services using the state machine approach: A tutorial. \emph{ACM Computing Surveys}, 22(4):299--319, 1990. \url{https://www.cs.cornell.edu/fbs/publications/SMSurvey.pdf}.

\bibitem{flp}
M. J. Fischer, N. A. Lynch, and M. S. Paterson. Impossibility of distributed consensus with one faulty process. \emph{Journal of the ACM}, 32(2):374--382, 1985. \url{https://groups.csail.mit.edu/tds/papers/Lynch/jacm85.pdf}.

\bibitem{bousquet}
G. D. Bousquet and J.-J. E. Slotine. A contraction based, singular perturbation approach to near-decomposability in complex systems. Research preprint, 2015. \href{https://arxiv.org/abs/1512.08464}{arXiv:1512.08464}.

\bibitem{dtn}
V. Cerf, S. Burleigh, A. Hooke, L. Torgerson, R. Durst, K. Scott, K. Fall, and H. Weiss. \emph{Delay-Tolerant Networking Architecture}. RFC 4838, Internet Research Task Force, April 2007. \url{https://www.rfc-editor.org/rfc/rfc4838}.

\bibitem{rao}
R. P. N. Rao, A. Stocco, M. Bryan, D. Sarma, T. M. Youngquist, J. Wu, and C. S. Prat. A direct brain-to-brain interface in humans. \emph{PLOS ONE}, 9(11):e111332, 2014. \href{https://doi.org/10.1371/journal.pone.0111332}{doi:10.1371/journal.pone.0111332}.

\bibitem{bennett}
C. H. Bennett. Notes on Landauer's principle, reversible computation, and Maxwell's demon. \emph{Studies in History and Philosophy of Modern Physics}, 34:501--510, 2003. \href{https://arxiv.org/abs/physics/0210005}{arXiv:physics/0210005}.

\bibitem{turing}
A. M. Turing. On computable numbers, with an application to the Entscheidungsproblem. \emph{Proceedings of the London Mathematical Society}, s2-42:230--265, 1936--1937. \url{https://www.cs.virginia.edu/~robins/Turing_Paper_1936.pdf}.

\bibitem{lei}
Z. Lei. On axially symmetric incompressible magnetohydrodynamics in three dimensions. Research paper, first posted 2012. \href{https://arxiv.org/abs/1212.5968}{arXiv:1212.5968}.
\end{thebibliography}
